\section{Облачный слой (Firebase Realtime Database)}

\subsection{Структура данных}

\begin{minted}{json}
{
  "current": { "t": 23.5, "h": 62.1, "g": 3200 },
  "history": {
    "1712700000": { "t": 22.1, "h": 64.2, "g": 2800 },
    "1712703600": { "t": 22.5, "h": 63.8, "g": 2900 }
  },
  "status": {
    "last_viewed": 1712703650
  }
}
\end{minted}

\subsection{Назначение веток в RTDB}

\begin{itemize}
    \item \texttt{/current} --- оперативное «live»-состояние для мгновенного отображения в UI.
    \item \texttt{/history/<hourEpoch>} --- архив телеметрии с часовым шагом.
    \item \texttt{/status/last\_viewed} --- технический маркер активности клиентского интерфейса.
\end{itemize}

В коде фронтенда предусмотрен парсинг и объекта, и JSON-строки, поэтому система устойчива к различным форматам payload в базе.

\subsection{Контракт данных}

\begin{itemize}
    \item \texttt{t}: температура, \texttt{number}, \textdegree C.
    \item \texttt{h}: влажность, \texttt{number}, \%.
    \item \texttt{g}: газ, \texttt{number} (сырой сигнал датчика, в UI приводится к \texttt{ppm} округлением).
    \item Все timestamp --- \textbf{Unix epoch, секунды, UTC}.
\end{itemize}

\subsection{Преимущества выбранной модели хранения}

\begin{enumerate}
    \item \textbf{Разделение горячих и холодных данных:} ветка \texttt{/current} не перегружает историю, а \texttt{/history} не мешает быстрому отображению.
    \item \textbf{Простота клиентской подписки:} для live-потока используется один наблюдаемый путь.
    \item \textbf{Предсказуемая ретроспектива:} исторические точки имеют фиксированную структуру ключей и легко сортируются.
\end{enumerate}

\section{Frontend (\texttt{iot-dashboard})}

\subsection{Стек}

\begin{itemize}
    \item React 19 + TypeScript.
    \item Vite.
    \item Material UI.
    \item Recharts.
    \item Firebase JS SDK (Realtime Database).
\end{itemize}

\subsection{Архитектурный принцип frontend}

Frontend построен по схеме «данные в хуках, отображение в компонентах». Это уменьшает связанность кода: бизнес-логика доступа к RTDB не смешивается с визуальным представлением.

\begin{figure}[htbp]
\centering
\begin{adjustbox}{max width=0.96\textwidth}
\begin{tikzpicture}[node distance=10mm and 12mm]
    \node[depbox] (app) {\texttt{App.tsx}};
    \node[depbox, below left=of app] (curc) {\texttt{components/}\\\texttt{CurrentData.tsx}};
    \node[depbox, below right=of app] (hisc) {\texttt{components/}\\\texttt{HistoryCharts.tsx}};

    \node[depbox, below left=of curc] (ucc) {\texttt{hooks/}\\\texttt{useCurrentData.ts}};
    \node[depbox, below right=of hisc] (uhc) {\texttt{hooks/}\\\texttt{useHistoryData.ts}};
    \node[depbox, below=of app] (tel) {\texttt{utils/telemetry.ts}};
    \node[depbox, below=of tel] (typ) {\texttt{types/index.ts}};
    \node[depbox, right=of tel] (fbc) {\texttt{config/firebase.ts}};

    \draw[deparrow] (app) -- (curc);
    \draw[deparrow] (app) -- (hisc);
    \draw[deparrow] (curc) -- (ucc);
    \draw[deparrow] (curc) -- (tel);
    \draw[deparrow] (hisc) -- (uhc);
    \draw[deparrow] (ucc) -- (fbc);
    \draw[deparrow] (uhc) -- (fbc);
    \draw[deparrow] (ucc) -- (typ);
    \draw[deparrow] (uhc) -- (typ);
    \draw[deparrow] (uhc) -- (tel);
\end{tikzpicture}

\end{adjustbox}
\caption{Диаграмма зависимостей модулей frontend}
\end{figure}

\subsection{Поток «живых» данных (\texttt{useCurrentData})}

\begin{enumerate}
    \item Создаётся подписка \texttt{onValue(ref(database, "current"))}.
    \item На каждом обновлении:
    \begin{itemize}
        \item payload проходит \texttt{parseTelemetryPayload};
        \item при успехе обновляется state \texttt{data}, \texttt{last\allowbreak Updated};
        \item при ошибке обновляется state \texttt{error}.
    \end{itemize}
    \item Параллельно работает heartbeat:
    \begin{itemize}
        \item сразу при монтировании;
        \item затем каждые \texttt{30 сек};
        \item в \texttt{/status/last\_viewed} пишется текущее UTC-время клиента.
    \end{itemize}
\end{enumerate}

Heartbeat --- это управляющий сигнал для прошивки. Благодаря ему MCU понимает, нужен ли пользователю live-поток прямо сейчас.

\subsection{Поток исторических данных (\texttt{useHistoryData})}

\begin{enumerate}
    \item По кнопке \texttt{Load data} компонент применяет диапазон дат.
    \item Хук делает \texttt{get(ref(database, "history"))}.
    \item Клиент фильтрует записи по диапазону \texttt{startTimestamp\allowbreak..\allowbreak endTimestamp}.
    \item Каждая запись:
    \begin{itemize}
        \item валидируется и парсится;
        \item преобразуется в \texttt{ChartDataPoint};
        \item сортируется по \texttt{timestamp}.
    \end{itemize}
    \item \texttt{HistoryCharts} строит график в выбранном режиме:
    \begin{itemize}
        \item все метрики;
        \item только температура;
        \item только влажность;
        \item только газ.
    \end{itemize}
\end{enumerate}

\subsection{UI-уровень}

\begin{itemize}
    \item \texttt{CurrentData}:
    \begin{itemize}
        \item 3 карточки (temperature/humidity/gas);
        \item индикатор режима Live;
        \item время последнего обновления.
    \end{itemize}
    \item \texttt{HistoryCharts}:
    \begin{itemize}
        \item выбор диапазона дат;
        \item валидация диапазона;
        \item интерактивная визуализация на Recharts.
    \end{itemize}
\end{itemize}

\subsection{Обработка ошибок и устойчивость UI}

\begin{itemize}
    \item Ошибки чтения не приводят к остановке приложения, а отображаются пользователю в виде контролируемого состояния.
    \item Невалидные payload не ломают компонент, а отсекаются на этапе парсинга.
    \item Отсутствие \texttt{/current} корректно интерпретируется как неактивный поток, а не как критическая ошибка.
\end{itemize}

\subsection{Вывод по облачному и клиентскому уровням}

Связка Firebase + React-хуки обеспечивает реактивную модель отображения с чётким разделением ответственности: RTDB хранит и синхронизирует состояние, hooks инкапсулируют доступ к данным, компоненты фокусируются на UX и визуальной достоверности представления.
